\chapter{Conclusion}
\label{chapter:conclusion}

The choice of QUIC implementation, or of QUIC versus TCP, has a large effect on
performance, but the best choice is not fixed: it depends on the network the
connection runs over and on what the deployment is optimising for. This thesis
set out to turn that decision into something measurable and automatic. It did so
by building a benchmarking framework, collecting a dataset across five
implementations and a wide range of network conditions, training a recommender
on it, and exposing that recommender through a live prober. The three research
questions from Chapter~\ref{chapter:introduction} can now be answered.

\textbf{RQ1: Does the best transport implementation depend on network conditions
and deployment priorities, or is there a single overall winner?} There is no
single winner, and the network conditions on their own are enough to move it.
Under clean traffic, TCP and msquic win most of the bandwidth/RTT grid, but
quiche takes over at high bandwidth-delay product, where its GSO-assisted send
path keeps the pipe full while the others cannot. Introduce 2\% packet loss and
quiche claims a large part of the grid from msquic, whose timeout risk climbs.
The RTT sweep shows the recommendation flipping from msquic at 10\,ms to TCP
beyond 60\,ms, and at the extreme corner of the grid (5\,Gbps, 300\,ms,
BDP\,$\approx$\,187\,MB) quiche and quic-go fall below 2\% link utilisation while
the link itself is idle. \emph{The winning implementation tracks the conditions
directly}: bandwidth, RTT through the BDP, and loss each move it. The deployment
profile then adds a second axis on top of this. For the very same condition the
datacenter profile favours quiche, which wins about 90\% of conditions under it,
while the lightweight profile favours TCP at over 99\%, and across the five
profiles the recommended implementation disagrees in up to 97.8\% of conditions.
Both the conditions and the priorities matter, which is exactly what makes a
single static default a poor choice.

\textbf{RQ2: Can a model trained offline on emulated measurements predict the
right implementation, and how accurately?} Yes, though the accuracy depends on
how cleanly the conditions separate the implementations. The recommender is a
two-stage random forest. Stage~1 predicts per-implementation timeout risk and is
near-perfect for four implementations, reaching 85.7\% for msquic, whose
failures sit on a genuinely noisy boundary at high BDP. Stage~2 ranks the
survivors for a given profile, scoring from 94.4\% on the datacenter profile
down to 67.8\% on the latency-sensitive one. That floor is not a modelling
failure: in 23.5\% of latency-sensitive conditions the top two implementations
score within 5\% of each other, so the label itself is close to a coin flip and
the model performs near the ceiling that ambiguity allows. Cross-validation and
out-of-bag scores agree closely across every profile, which means the
condition-level GroupKFold splits are not leaking repetitions and the reported
numbers are honest.

\textbf{RQ3: Can the recommender run against a live connection?} Yes. The prober
was run inside the same Docker topology used for data collection. It reproduced
the configured RTT to within 1\,ms, recovered the other conditions (with loss
correctly measured higher than configured, because each packet crosses the
impairment twice), and produced a recommendation with a confidence value and a
STABLE/UNSTABLE flag on every profile tested. It also behaves sensibly at the
edges. In the deep-space case it eliminates three QUIC stacks on timeout risk
and recommends TCP at full confidence. In the mobile-lossy case, where measured
loss ran above the training range, Stage~2 rates quiche the best raw performer
but Stage~1 eliminates it on timeout risk, so the recommender falls back to TCP
and reports that quiche was the faster but less reliable option.

% =============================================================================
\section{Limitations}
\label{sec:conclusion-limitations}

These results come with limits worth stating plainly. The measurements are from
controlled \texttt{tc/netem} emulation on a single Linux host, not the open
internet, and they cover single-stream bulk transfer rather than the
multiplexed, request-response traffic of a real web workload. The five
implementations are tested as their documentation configures them for
throughput, which means the comparison is deliberately not apples-to-apples:
differences that come from one project shipping an optimisation another lacks
are kept rather than normalised away, because that is what a developer choosing
between them would actually experience. The msquic flow-control boundary that
the \texttt{fc\_limited} feature encodes is an empirical, configuration-dependent
threshold, the 32\,MB initial per-stream window we set, not a hard protocol
limit, and a different window setting could move it. Two of the five
implementations, mvfst and quic-go, win so rarely under every profile that they
are thinly represented as Stage~2 labels, so the recommender's ability to choose
them confidently is correspondingly weak. None of these undermine the central
result, but they bound how far it should be generalised.

Within those bounds, the thesis shows that picking a transport implementation is
a decision worth making per deployment rather than once, and that a small,
explainable model trained on emulated measurements can make it from conditions
that are cheap to observe.
